% ============================================================
% pure 报告 — 规范模板（xelatex + ctex）
% 主题: quda（QUDA 1.1.0）核心部分穷尽剖析
% 编译: xelatex 两遍; 交付前 Overfull / Float too large 均为 0
% ============================================================
\documentclass[UTF8]{ctexart}
\usepackage[margin=2.5cm]{geometry}
\usepackage{graphicx}
\usepackage{amsmath}
\usepackage{amssymb}
\usepackage{microtype}
\usepackage{listings}
\usepackage{xcolor}
\usepackage{booktabs}
\usepackage{longtable}
\usepackage{xltabular}
\usepackage{tabularx}
\usepackage{hyperref}
\usepackage{fancyhdr}
\usepackage[most]{tcolorbox}
\usepackage{titlesec}

\definecolor{accent}{HTML}{1F4E79}
\definecolor{evidence}{HTML}{1E8449}
\definecolor{reference}{HTML}{8C6D1F}
\definecolor{conclusion}{HTML}{8B1A1A}
\definecolor{codebg}{HTML}{F4F6F7}
\definecolor{algo}{HTML}{E67E22}
\definecolor{phys}{HTML}{6C3483}
\definecolor{map}{HTML}{C2185B}
\definecolor{warn}{HTML}{D35400}
\definecolor{hl}{HTML}{FDF6E3}

\setlength{\emergencystretch}{3em}
\hypersetup{colorlinks=true,linkcolor=accent,urlcolor=map}

\titleformat{\section}{\Large\bfseries\color{accent}}{\thesection}{1em}{}
\titleformat{\subsection}{\large\bfseries\color{phys}}{\thesubsection}{1em}{}
\titleformat{\subsubsection}{\normalsize\bfseries\color{phys!70!black}}{\thesubsubsection}{1em}{}

\newtcolorbox{conclbox}{breakable,colback=conclusion!6,colframe=conclusion,title=结论}
\newtcolorbox{refbox}{breakable,colback=reference!6,colframe=reference,title=参考背景}
\newtcolorbox{warnbox}{breakable,colback=warn!6,colframe=warn,title=局限与风险}
\newtcolorbox{keybox}{breakable,colback=hl,colframe=accent,title=要点}

\lstset{
  basicstyle=\ttfamily\footnotesize,
  backgroundcolor=\color{codebg},
  frame=single,
  language=C++,
  firstnumber=auto,
  keywordstyle=\color{accent}\bfseries,
  commentstyle=\color{evidence!70!black},
  stringstyle=\color{phys},
  breaklines=true,
  breakatwhitespace=false,
  postbreak=\mbox{\textcolor{accent}{$\hookrightarrow$}\space},
  keepspaces=true,
  columns=flexible,
  numbers=left,numberstyle=\tiny\color{phys!60!black},
}

\title{quda（QUDA 1.1.0）核心部分穷尽剖析报告}
\author{opencode pure}
\date{2026-08-17}

\begin{document}
\maketitle

\begin{abstract}
本报告对 \url{/root/PyQCU/refer/git-rep/quda}（QUDA 1.1.0）的核心部分进行穷尽剖析。
项目性质判定为\textbf{代码-物理交叉向}：完整 GPU 格点 QCD 库，物理算子（Dirac/dslash/
multigrid/特征求解）为核心，性能工程（模板内核/JIT/后端抽象）为支撑。穷尽枚举 12 个核心
候选（C1--C12），其中\textbf{深挖 7 / 浅析 3 / 排除 2}：深挖部分含物理图像、公式推导链、
代码-物理映射表与证据。独特竞争力：自适应 multigrid、模板化 dslash 内核族、TRLM 特征求解、
SU(3) 重构压缩、多后端（CUDA/HIP）抽象。非独立 git 仓库。
\end{abstract}

\tableofcontents

\section{项目定位与核心部分清单}

\subsection{项目性质判定}
\begin{keybox}
\textbf{性质判定：交叉向。}物理对象密集（Dirac 算子族、规范场、clover、色旋量场、multigrid、
特征求解），算法/内核为实现核心。判定依据：\texttt{include/dirac\_quda.h}（12 个 Dirac 类）、
\nolinkurl{include/multigrid.h:263}（MG 类）、\texttt{lib/multigrid.cpp:1131}（V-cycle）。
\end{keybox}

\subsection{核心部分总清单（穷尽枚举）}
\begin{xltabular}{\textwidth}{lXXX}
\toprule
\textcolor{map}{编号} & 候选 & 三态 & 理由 \\
\midrule
\endhead
C1 & 自适应多重网格（MG）& 深挖 & 求解性能核心 \\
C2 & 粗格算子与转移 & 深挖 & MG 正确性核心 \\
C3 & Wilson dslash 内核 & 深挖 & 最常用算子内核 \\
C4 & 域墙/Möbius 算子 & 深挖 & 独特作用量族 \\
C5 & Krylov 求解器族 & 深挖 & 反演核心 \\
C6 & TRLM 特征求解 & 深挖 & 低模物理核心 \\
C7 & SU(3) 重构压缩 & 深挖 & 存储/带宽优化核心 \\
C8 & Clover 构造与力 & 浅析 & 支撑性组件 \\
C9 & HMC 力（gauge/clover/HISQ）& 浅析 & 经 C API 暴露 \\
C10 & halo 通信与 GPU-direct & 浅析 & 性能优化 \\
C11 & MMA 粗算子（张量核）& 排除 & 可选特性 \\
C12 & 距离预条件 & 排除 & 可选特性 \\
\bottomrule
\caption{核心部分清单三态（数据来源: 全库索引实测）}
\end{xltabular}

\section{核心部分深度剖析}

\subsection{C1 自适应多重网格（MG）}
\textbf{物理图像：}格点 Dirac 算子近零模导致 Krylov 求解器收敛退化，MG 以聚合构造粗格
捕捉低频。\textbf{机制：}\texttt{lib/multigrid.cpp:1131}（MG::operator() V-cycle）、
\texttt{:1275}（\texttt{generateNullVectors}）、\texttt{:1646}
（\texttt{generateEigenVectors}）、\texttt{multigrid.h:263}（MG 类）。
\textbf{推导链：}null 向量经平滑-正交化生成 $\rightarrow$ 构造 $P$ $\rightarrow$ 粗算子
$D_c=P^\dagger D P$（粗层 DiracCoarse，\texttt{multigrid.h:73}）。
\textbf{为什么自适应：}物理质量参数近临界时近零模密集，自适应 null 向量生成保证 $P$
逼近低模空间。

\subsection{C2 粗格算子与转移}
\textbf{物理图像：}粗细格转移（prolongator/restrictor）与粗算子 $D_c$ 是 MG 正确性关键。
\textbf{机制：}\nolinkurl{include/kernels/coarse_op_kernel.cuh}、
\nolinkurl{include/kernels/restrictor.cuh}、\texttt{prolongator.cuh}、
\texttt{lib/coarse\_op.in.cu}、\texttt{lib/restrictor.in.cu}、\texttt{prolongator.in.cu}。
\textbf{推导链：}
\begin{equation}\label{eq:mggalerkin}
D_c = P^\dagger D P,\quad \psi_{\text{fine}} = P\,\psi_{\text{coarse}},\quad \eta_{\text{coarse}} = P^\dagger \eta_{\text{fine}}.
\end{equation}
\textbf{极限校验：}$P^\dagger P = I$（正交）时粗算子谱保留；MMA 版（\nolinkurl{coarse_op_kernel_mma.cuh}）
用张量核加速。

\subsection{C3 Wilson dslash 内核}
\textbf{物理图像：}Wilson 算子是格点 QCD 基础运算。\textbf{机制：}
\nolinkurl{include/kernels/dslash_wilson.cuh:18}（WilsonArg）、\texttt{:117}（apply 内核）、
\texttt{include/dslash.h:33}（Dslash 启动模板）、\nolinkurl{include/instantiate_dslash.h}
（按精度/重构/颜色实例化）、\nolinkurl{lib/dslash_wilson.in.cu}。
\textbf{推导链：}
\begin{equation}\label{eq:dslash}
D\psi = \sum_\mu\left[U_\mu(x)(1-\gamma_\mu)\psi(x+\hat\mu) + U_\mu^\dagger(x-\hat\mu)(1+\gamma_\mu)\psi(x-\hat\mu)\right].
\end{equation}
\textbf{为什么模板实例化：}精度（half/single/double）×重构（18/12/8/13）组合爆炸，
configure\_file 批量生成内核实例（\texttt{lib/CMakeLists.txt}）。

\subsection{C4 域墙/Möbius 算子}
\textbf{物理图像：}域墙（Domain Wall）/Möbius 作用量引入第 5 维 $L_s$，恢复手征对称性。
\textbf{机制：}\nolinkurl{include/dirac_quda.h:851}（DiracDomainWall）、\texttt{:973}
（DiracMobius）、\texttt{:1073}（DiracMobiusEofa）、\nolinkurl{lib/dirac_domain_wall.cpp}、
\texttt{lib/dirac\_mobius.cpp}、\nolinkurl{lib/dslash_domain_wall_*.in.cu}。
\textbf{推导链：}$D_{DW}$ 以 Wilson 核沿第 5 维耦合（M5 负质量参数）。
\textbf{为什么 EOFA：}Möbius EOFA 变体消除部分怪级部分子传播子，优化反演。

\subsection{C5 Krylov 求解器族}
\textbf{物理图像：}狄拉克反演 $D\psi=b$ 的迭代求解。\textbf{机制：}
\nolinkurl{include/invert_quda.h:381}（Solver 基类）、CG:715、CGNE:781、CGNR:816、PCG:886、
BiCGstab:949、BiCGstabL:992、GCR:1105、MR:1172、CACG:1211、CAGCR:1284、SD:1334、
GMRESDR:1629；\texttt{lib/inv\_cg\_quda.cpp:63}、\nolinkurl{inv_bicgstab_quda.cpp:79}、
\nolinkurl{inv_bicgstabl_quda.cpp:394} 等。
\textbf{为什么多样：}不同作用量（非 Hermite/带质量偏移/多层）需不同收敛策略；混合精度
CG 加速（精细/粗精度双跑）。

\subsection{C6 TRLM 特征求解}
\textbf{物理图像：}低模特征值对物理量（传播子、topological susceptibility）关键。
\textbf{机制：}\nolinkurl{include/eigensolve_quda.h:16}（EigenSolver）、\texttt{:336}
（TRLM）、\texttt{:457}（TRLM3D）、\texttt{:567}（IRAM）；\texttt{lib/eig\_trlm.cpp:39}、
\texttt{lib/eig\_iram.cpp}。\textbf{为什么 TRLM：}厚重启 Lanczos 适合大规模稀疏非 Hermite
算子，配合 deflation（\texttt{include/deflation.h:75}）加速反演。

\subsection{C7 SU(3) 重构压缩}
\textbf{物理图像：}规范 link $U_\mu(x)\in SU(3)$ 有 18 实分量，可用 12/8/13 分量重构压缩
（重构 18：$\{u,v\}$ 两列 12 元 + 第三列由行列式西性重构）。\textbf{机制：}
\nolinkurl{include/gauge_field_order.h}（重构 8/12/13/18 元素压缩）、
\nolinkurl{include/su3_project.cuh}。\textbf{为什么：}内存带宽是 dslash 性能瓶颈，压缩
降低显存访问量。\textbf{极限校验：}重构 18（不压缩）为基准；重构 8 最激进但需额外
平方根计算。

\section{独特性与竞争力评估}
\begin{xltabular}{\textwidth}{lXX}
\toprule
\textcolor{algo}{核心} & 独特竞争力 & 对比朴素实现 \\
\midrule
\endhead
自适应 MG + 特征向量 & 格子无关收敛 & 纯 Krylov \\
模板化 dslash 内核族 & 精度×重构组合实例化 & 单一内核 \\
TRLM/IRAM 特征求解 & 大规模低模求解 & 幂法 \\
SU(3) 重构压缩 & 带宽优化 & 全 18 分量 \\
多后端抽象（CUDA/HIP）& 跨硬件 & 平台绑定 \\
\bottomrule
\caption{独特性评估（数据来源: 源码实测）}
\end{xltabular}

\section{关键片段与公式}
\begin{lstlisting}[language=C++]
void MG::operator()(cvector_ref<ColorSpinorField> &x,
                    cvector_ref<const ColorSpinorField> &b)
\end{lstlisting}
参考源：\texttt{lib/multigrid.cpp:1131}（MG V-cycle 入口）

核心公式汇总：
\begin{align}
D_c &= P^\dagger D P \quad(\text{Galerkin}), \label{eq:q1}\\
D &= \sum_\mu\left[U_\mu(1-\gamma_\mu)\delta_{+} + U_\mu^\dagger(1+\gamma_\mu)\delta_{-}\right], \label{eq:q2}\\
\text{重构 12} &: U = \{u,v\},\ w = u\times v^{\dagger}/(\det u \det v). \label{eq:q3}
\end{align}

\section{参考源清单}
{\small
\begin{xltabular}{\textwidth}{XXX}
\toprule
\textcolor{evidence}{参考源} & 类型 & 内容/作用 \\
\midrule
\endhead
\texttt{\nolinkurl{include/multigrid.h:263}} & 仓库 & MG 类 \\
\texttt{\nolinkurl{lib/multigrid.cpp:1131}} & 仓库 & MG::operator() \\
\texttt{\nolinkurl{lib/multigrid.cpp:1275}} & 仓库 & generateNullVectors \\
\texttt{\nolinkurl{lib/multigrid.cpp:1646}} & 仓库 & generateEigenVectors \\
\texttt{\nolinkurl{include/kernels/coarse_op_kernel.cuh}} & 仓库 & 粗算子内核 \\
\texttt{\nolinkurl{include/kernels/dslash_wilson.cuh:18,117}} & 仓库 & WilsonArg/apply \\
\texttt{\nolinkurl{include/dslash.h:33}} & 仓库 & Dslash 启动模板 \\
\texttt{\nolinkurl{lib/dslash_wilson.in.cu}} & 仓库 & 内核模板实例化 \\
\texttt{\nolinkurl{include/dirac_quda.h:851,973,1073}} & 仓库 & DW/Möbius/EOFA \\
\texttt{\nolinkurl{include/invert_quda.h:381,715,949,1629}} & 仓库 & Solver/CG/BiCGstab/GMRESDR \\
\texttt{\nolinkurl{lib/inv_cg_quda.cpp:63}} & 仓库 & CG 实现 \\
\texttt{\nolinkurl{lib/inv_bicgstab_quda.cpp:79}} & 仓库 & BiCGstab 实现 \\
\texttt{\nolinkurl{include/eigensolve_quda.h:16,336,567}} & 仓库 & EigenSolver/TRLM/IRAM \\
\texttt{\nolinkurl{lib/eig_trlm.cpp:39}} & 仓库 & TRLM 实现 \\
\texttt{\nolinkurl{include/deflation.h:75}} & 仓库 & Deflation \\
\texttt{\nolinkurl{include/gauge_field_order.h}} & 仓库 & SU(3) 重构压缩 \\
\texttt{\nolinkurl{lib/color_spinor_field.cpp:1256}} & 仓库 & exchangeGhost \\
\texttt{\nolinkurl{lib/communicator_mpi.cpp:106}} & 仓库 & comm\_init \\
\texttt{\nolinkurl{CMakeLists.txt:77,266}} & 仓库 & 目标类型/MG 开关 \\
\texttt{\nolinkurl{lib/targets/cuda/target_cuda.cmake:12}} & 仓库 & GPU 架构默认 \\
\bottomrule
\end{xltabular}
}

\section{结论}
\begin{conclbox}
QUDA 的核心竞争力是\textbf{自适应 MG + 模板化内核 + 特征求解 + 重构压缩 + 多后端抽象}的组合：
以 configure\_file 批量实例化应对精度×重构组合爆炸，以 SU(3) 重构优化带宽，以 TRLM/deflation
支撑低模物理，是生产级 GPU 格点 QCD 库的范式。穷尽枚举 12 个候选，深挖 7 项均有物理图像
与代码-物理映射证据。
\end{conclbox}
\begin{warnbox}
\textcolor{warn}{局限：}未实测运行（无构建产物）；构建选项繁多学习成本高；本版本无 HMC
集成器（力经 C API 暴露）；MMA 与距离预条件为可选特性未纳入深挖；非独立 git 仓库。
\end{warnbox}

\end{document}